% Proposed edition Chapter 281. Source-bound directory, not a replacement
% for the separately maintained presentations or their proof obligations.
\chapter{Directory of Logical Systems}
\label{sec:p14-directory-of-logical-systems}
\index{logical system!directory}\index{consequence relation!directory}

The name of a logic is not enough to identify its consequence relation.
The language, admissible structures, evaluation points, designated values,
premise convention and proof calculus must also agree. This directory
locates those choices in the handbook. Its entries preserve the scope of
the linked presentation: a propositional result does not classify a
quantified extension, a theoremhood result does not establish strong
completeness, and a rule sheet does not supply an omitted adequacy proof.
The substantive treatments retain their own definitions, examples,
sources and result standing \cite{Enderton2001,Manzano1996,
TroelstraSchwichtenberg2000,BlackburnDeRijkeVenema2001}.

\section{Languages, semantic structures, and consequence relations}
\label{sec:p14-directory-interfaces}

\paragraph{What an entry identifies.}
A language specifies its primitive expressions, binding operations and
formation rules. A model class specifies admissible interpretations, not
merely a vocabulary of possible interpretations. Satisfaction may be
evaluated at a valuation, a structure and assignment, a world, a path
position, an information state or a probability structure. These coordinates
belong to the statement of consequence. In the familiar truth-preserving
case, semantic consequence says that every admissible evaluation satisfying
all premises satisfies the conclusion. Derivability instead refers to
finite proofs in a specified calculus. Their agreement is an adequacy
result, not part of the definition of either relation. The foundational
satisfaction and consequence conventions are in
\cref{sec:prop-syntax-semantics-02,sec:fol-semantics-02}; the
intuitionistic and support-based relations are fixed separately in
\cref{def:p08-ipl-consequence,def:p08-inq-general-entailment}.

\begin{longtable}{@{}p{.20\textwidth}p{.35\textwidth}p{.37\textwidth}@{}}
\caption{Evaluation and premise interfaces used by the directory.}
\label{tab:p14-directory-consequence}\\
\toprule
Interface & Semantic data and conclusion & What must not be silently changed\\
\midrule
\endfirsthead
\toprule
Interface & Semantic data and conclusion & What must not be silently changed\\
\midrule
\endhead
\bottomrule
\endfoot
Classical truth preservation & A valuation, or a structure with an
assignment, makes each premise and the conclusion true. & Nonempty domains,
equality, the treatment of open formulas and the chosen signature; see
\cref{sec:fol-semantics-02}.\\\addlinespace
Designated-value preservation & A matrix evaluation designates every
premise only if it designates the conclusion. & The designated set is part
of the logic. K3 and LP use the same operations with different designation;
see \cref{entry:p08-k3,entry:p08-lp}.\\\addlinespace
Local modal consequence & Premises and conclusion are evaluated at one
common world of each admissible model. & Global premises hold at every
world and define a different relation. Necessitation is theorem-only in
the local premise calculi; see \hyperref[sheet:hilbert-k]{the K sheet}.\\\addlinespace
Ordered or multiset sequent validity & An antecedent product is below the
succedent interpretation in the declared algebra. & Order, multiplicity,
empty contexts, units and structural rules; see
\cref{entry:p08-imll,def:p08-lambek-semantics}.\\\addlinespace
Support consequence & Every information state supporting all premises
supports the conclusion. & Support for a question settles it; this is not
truth at a single world or a claim about an agent's investigation;
see \cref{def:p08-inq-general-entailment}.\\\addlinespace
Defeasible selection & A declared preference, extension, expansion or
minimal-model policy selects the evaluations used for the conclusion. &
Smoothness, ranking, grounding, existence and the some/all choice;
see \cref{def:p08-preferential-model,entry:p08-reiter-default}.\\\addlinespace
Program correctness & An assertion constrains executions or a heap
interpretation of a command. & Partial correctness permits divergence;
total correctness adds termination. Heap splitting and locality are
separate hypotheses; see \cref{def:p09-hoare-triple,def:p09-separation-logic}.\\
\end{longtable}

\paragraph{Theoremhood, finite consequence and arbitrary premises.}
Validity is truth without premises. Weak completeness concerns valid
formulas; strong completeness concerns the specified premise sets.
Finite-premise completeness does not imply arbitrary-premise completeness.
The distinction is decisive for full second-order semantics, LTL, CTL*,
common knowledge and several standard real-valued chains. Likewise,
multiple-conclusion matrix consequence requires at least one designated
conclusion, whereas the single-conclusion relation fixes one conclusion.
The K3--LP duality in \cref{thm:p08-k3-lp-facts} explicitly uses finite
sets on both sides; it is not an identity of their single-conclusion logics.

\paragraph{Logic, theory, calculus and reasoning method.}
A theory adds nonlogical commitments to a language and consequence relation.
Arithmetic and set-theoretic axiom systems therefore retain their theory
identities rather than becoming interchangeable names for first-order logic.
A calculus is a proof presentation; natural deduction, sequents, Hilbert
proofs, tableaux and resolution can serve the same classical consequence
relation while retaining different rules and side conditions. A normative
argument, a statistical procedure, a causal model or an empirical model of
human reasoning does not acquire a logical language and adequacy theorem
merely by appearing in a handbook of reasoning. Those methods are located
by \cref{sec:p14-reasoning-task-finder,sec:p14-reasoning-methods-and-evidence};
this directory includes them only when a particular formal language and
inferential interpretation are specified.

\section{Classical, constructive, modal, and many-valued systems}
\label{sec:p14-directory-classical-constructive}

\paragraph{Classical propositional logic.}
\label{entry:p14-directory-cpl}
The inherited language has propositional variables, falsum and the
primitive classical Boolean connectives. Its structures are two-valued
valuations and its semantic consequence is truth preservation over all
valuations. The syntax and recursive clauses are in
\cref{sec:prop-syntax-semantics-01,sec:prop-syntax-semantics-02}; the
natural-deduction and sequent treatments and the separately attributed
Hilbert/tableau comparison retain their distinct proof objects. Soundness,
completeness and compactness are located in
\cref{sec:prop-metatheory-01,sec:prop-metatheory-03}.
Finite truth tables decide consequence; the finite-input decision and
encoding distinctions are collected in
\cref{sec:p14-decidability-and-complexity}. Equality is not a primitive
of this propositional language \cite{Enderton2001,Mendelson1997}.

\paragraph{Classical first-order logic and strict many-sorted logic.}
\label{entry:p14-directory-fol}
First-order formulas use the declared relation, function and constant
symbols, object variables, Boolean connectives and object quantifiers.
Structures have nonempty domains, total functions and relation extensions;
assignments interpret free variables. Equality and the equality-free
presentations have distinct rule surfaces, in particular
\hyperref[sheet:lk-eq]{LK with equality} and
\hyperref[sheet:hilbert-fol]{the equality-free Hilbert presentation}.
Use \cref{def:p02-fol-signature,sec:fol-semantics-02} for the semantic
interface and \cref{sec:fol-metatheory-03,sec:fol-metatheory-04} for completeness and
compactness. These unrestricted-model results are not finite-model
compactness or a decision procedure for every first-order input.
The strict many-sorted variant in \cref{sec:p02-many-sorted} instead
has disjoint nonempty sort domains, typed functions and relations, and
no cross-sort equality. Its declared countable signature and relativization
background are part of the reduction to one-sorted logic; allowing empty or
overlapping sorts would change that presentation \cite{Enderton2001,Manzano1996}.

\paragraph{Categorical syllogistic fragments.}
\label{entry:p14-directory-syllogistic}
The modern systems \(\mathcal S\) and \(\mathcal S^\dagger\) in
\cref{sec:p02-syllogistic-inference} reason with quantified term formulas
over unary predicates in a nonempty universe. The latter admits the source's
additional literal-subject forms; neither is unrestricted Boolean
first-order syntax. Predicate extensions may be empty. Their direct
calculus, model construction and finite decision method are in that
treatment, rather than supplied by the directory. Universal inclusion
does not entail an existential witness without an existence premise.
These are the separately attributed 2008 preprint fragments, not an
unmodified presentation of Aristotle's historical syllogistic
\cite{PrattHartmannMoss2008SyllogisticV1}.

\paragraph{Relational second-order logic: full and Henkin regimes.}
\label{entry:p14-directory-sol}
The relational language quantifies over objects and relations of each
declared arity. Under full semantics, a relation quantifier ranges over
every relation of that arity on the object domain; under relational Henkin
semantics it ranges over a supplied relation domain. Full semantics and
its noncompactness and effective-completeness limitations are in
\cref{sec:second-order-logic-02,thm:p02-sol-full-noncompact,
thm:p02-sol-full-no-effective-completeness}. An arbitrary Henkin
structure is not automatically closed under comprehension or choice.
The distinct matched class and its consequence relation are in
\cref{def:p02-sol-matched-henkin,def:p02-sol-henkin-consequence}.
\Cref{thm:p02-sol-derived-adequacy} supplies adequacy for the
translation-defined derivability relation under its ambient choice
assumption; \cref{def:p02-sol-native-boundary} does not assert completeness
of the native displayed rule list. Neither result provides an effective
complete calculus for full second-order validity \cite{Henkin1950,Manzano1996}.

\paragraph{Infinitary and generalized-quantifier languages.}
\label{entry:p14-directory-infinitary}
The displayed \(L_{\omega_1\omega}\) grammar admits countable Boolean
combinations with a finite total free-variable set, but only finite
quantifier strings. Its consequence uses the extended structure/assignment
satisfaction clauses in \cref{def:p02-lomega-one-omega}.
The noncompactness and countable-model Scott-sentence statements are at
\cref{thm:p02-lomega-one-omega-calibration}; the cited Karp completeness
boundary names genuinely well-founded infinitary proof trees, not finite
proofs. Generalized quantifiers instead add a specified isomorphism-invariant
class of domain/relation pairs to the satisfaction clauses.
\Cref{def:p02-generalized-quantifier,sec:p02-infinitary-generalized}
keep that class as an explicit parameter. The compactness or
axiomatizability of one selected quantifier is not a family-wide result
\cite{Karp1964,Scott1965,Mostowski1957,Lindstrom1966,Keisler1970}.

\paragraph{Intuitionistic propositional and first-order logic.}
\label{entry:p14-directory-intuitionistic}
IPL, also called IPC in the literature, uses falsum, conjunction,
disjunction and implication, with negation defined through falsum.
Its Kripke models use persistent atomic valuations, and implication tests
all later worlds. Local consequence preserves forcing at one common
world. The exact partially ordered presentation, its complete adequacy
argument and finite filtration are at
\cref{entry:p08-ipl,sec:p08-ipl-kripke-completeness}; the separately stated
preorder semantics is at \cref{def:p08-ipl-kripke-structure}.
Excluded middle and double-negation elimination are not valid in general.
First-order intuitionistic logic adds the domain and assignment conventions
of \cref{def:p08-iqc-expanding-domain}; a propositional finite-model argument does not
give its decision classification. \hyperref[sheet:lj]{LJ} and
\hyperref[sheet:nj]{NJ} retain their own single-succedent and discharge
conditions \cite{Kripke1965,TroelstraVanDalen1988,TroelstraSchwichtenberg2000}.

\paragraph{Intermediate logics KC and LC.}
\label{entry:p14-directory-intermediate}
These extend the IPL theorem set in its propositional language. KC adds
weak excluded middle and has the stated finite rooted-frame profile with
a greatest world; LC, or G\"odel--Dummett logic, adds prelinearity and
is characterized by linearly ordered intuitionistic frames.
\Cref{thm:p08-kc-profile,thm:p08-lc-profile,tab:p08-intermediate-logics}
record their finite-model and disjunction-property boundaries.
The explicit modern Hilbert presentations and their strong local
frame proofs are \cref{thm:p08-mi-intermediate-adequacy};
\cref{thm:p08-mi-finite-decision} supplies terminating finite-input
decisions while retaining the original-source distinctions.
Their strict chain between IPC and CPC is only a subdiagram of the full
lattice of intermediate logics. A theorem-set extension is not a claim
that every arbitrary Kripke frame class determines one of these two systems
\cite{Jankov1968WeakExcludedMiddle,Dummett1959,ChagrovZakharyaschev1997}.

\paragraph{Typed calculi and constructive interpretations.}
\label{entry:p14-directory-typed}
The arrow-only simply typed lambda calculus relates annotated terms,
simple types and typing contexts; normalization and subject reduction
belong to that judgmental system. Church-style System F adds type
abstraction, type application and universal types. Its explicitly annotated
checking problem is not Curry-style erased-term typability. The exact
languages, reductions and type judgments are in
\cref{sec:p07-lambda-calculi-and-type-theories}; the separate System F
normalization and confluence/checking developments are linked from
\cref{thm:p07-system-f-strong-normalization,thm:p07-system-f-confluence,
thm:p07-system-f-church-checking}. These are typing and reduction judgments,
not an unspecified truth-preservation semantics. The dependent-type entry
is the Abel--Coquand--Dybjer core with dependent functions and one
Russell-style universe. Its typed equality result does not cover all
Martin--L\"of extensions, inductive types, univalence or cubical computation;
see \cref{thm:p07-mltt-normalization-decision,def:p07-hott-cubical-boundary}
\cite{GirardLafontTaylor1989,AbelCoquandDybjer2007}.

\paragraph{Normal propositional modal systems.}
\label{entry:p14-directory-normal-modal}
All rows below use the classical propositional modal language with box
and its defined diamond. Consequence is local truth preservation in the
specified Kripke class. The Hilbert premise convention keeps necessitation
theorem-only. The cards, correspondence results and exact sheets are
maintained in Part VIII, not replaced by the frame descriptions here
\cite{Chellas1980,HughesCresswell1996,BlackburnDeRijkeVenema2001,
ChagrovZakharyaschev1997}.

\begin{longtable}{@{}p{.13\textwidth}p{.42\textwidth}p{.37\textwidth}@{}}
\caption{The named normal-modal catalog and its actual destinations.}
\label{tab:p14-directory-normal-modal}\\
\toprule
System & Admissible frames or validity regime & System treatment and distinction\\
\midrule
\endfirsthead
\toprule
System & Admissible frames or validity regime & System treatment and distinction\\
\midrule
\endhead
\bottomrule
\endfoot
\(K\) & All nonempty Kripke frames. & \cref{entry:p08-modal-k}; no T,
4 or 5 condition is automatic.\\\addlinespace
\(KD\) & Serial frames. & \cref{entry:p08-modal-d}; the alias D is
source-bound, and seriality is not reflexivity.\\\addlinespace
\(KT\) & Reflexive frames. & \cref{entry:p08-modal-t}; also called T
in the specified source convention.\\\addlinespace
\(KTB\) & Reflexive symmetric frames. & \cref{entry:p08-modal-b};
Hughes--Cresswell's B is not Chellas's symmetry-only KB.\\\addlinespace
\(K4\) & Transitive frames. & \cref{entry:p08-modal-k4}; reflexivity
is not included.\\\addlinespace
\(S4\) & Reflexive transitive frames. & \cref{entry:p08-modal-s4};
also KT4.\\\addlinespace
\(KD45\) & Serial transitive Euclidean frames. &
\cref{entry:p08-modal-kd45}; not a factive knowledge system.\\\addlinespace
\(S4.2\) & Reflexive transitive frames directed above each world. &
\cref{entry:p08-modal-s42}; directed convergence is an extra condition.\\\addlinespace
\(S4.3\) & Reflexive transitive frames weakly connected above each world. &
\cref{entry:p08-modal-s43}; the exact weak-connectedness convention is
in \cref{def:p08-modal-frame-properties}.\\\addlinespace
\(S5\) & Equivalence frames. & \cref{entry:p08-modal-s5}; the
registered presentation uses T and 5.\\\addlinespace
GL & Formula validity on finite transitive irreflexive frames. &
\cref{entry:p08-modal-gl}; the finite-frame theoremhood result is not
arbitrary-premise strong completeness.\\\addlinespace
Grz & Formula validity on finite partial orders. &
\cref{entry:p08-modal-grz}; a separate branch, not a modal-cube alias.\\
\end{longtable}

\paragraph{Neighborhood modal systems.}
\label{entry:p14-directory-neighborhood}
Chellas's classical systems use a neighborhood family at each world;
box truth means that the truth set of its argument belongs to that family.
E imposes no closure condition. EM adds upward closure, EC binary
intersection closure, EN the whole-world-set neighborhood, and EMC the
first two conditions together. EMCN is the normal K boundary.
\Cref{tab:p08-neighborhood-systems,sec:p08-n-b-006-03}
give the separate names, source language with primitive possibility,
definitional bridge, local consequence and canonical constructions.
Chellas's system-name R for regular logic belongs to this neighborhood
presentation; it is not relevant implication \(\mathbf R\).
The original Scott and Montague frameworks and Pacuit's later treatment
retain separate source identities \cite{Scott1970,Montague1970Pragmatics,Chellas1980,Pacuit2017}.

\paragraph{K3, LP, FDE and weak Kleene.}
\label{entry:p14-directory-finite-matrices}
The K3, LP and FDE entries fix the constant-free language of negation,
conjunction and disjunction. K3 uses the strong Kleene three-valued
operations with only true designated; LP designates both true and the
middle, both-supported value. FDE uses four support pairs, separating
neither support from both support and designating exactly positive-support
values. Consequence preserves designation, not equality of values.
\Cref{entry:p08-k3,entry:p08-lp,entry:p08-fde,def:p08-fde-matrix}
give their complete matrix interfaces and the finite-matrix tableau route.
K3 has no tautologies in this fixed language; LP has the classical
tautologies but not classical consequence; FDE is both nonexplosive and
paracomplete. The weak Kleene matrix in
\cref{def:p08-weak-kleene-matrix} instead makes the middle value
infectious. It is not a reconstruction of Bochvar's richer external
language \cite{Kleene1952,Priest1979,Dunn1976,Belnap1977,OmoriWansing2017}.

\paragraph{Finite and infinite \L{}ukasiewicz systems.}
\label{entry:p14-directory-lukasiewicz}
For each fixed number of values at least two, the finite chain uses
implication and negation with the displayed MV operations and only value
one designated. The infinite standard semantics uses the real unit
interval. The language, operations and distinct calculi are in
\cref{def:p08-lukasiewicz-matrices,entry:p08-lukasiewicz-finite,
entry:p08-lukasiewicz-infinite}. Exhaustive finite-matrix evaluation
does not decide the interval semantics by the same procedure. Standard
theoremhood and finite-premise completeness do not become arbitrary-premise
standard completeness; the all-MV-algebra regime is separate. McNaughton's
truth-function representation is located at the source-bound infinite
entry \cite{Lukasiewicz1920,LukasiewiczTarski1930,McNaughton1951,
Hajek1998,CignoliDOttavianoMundici2000}.

\paragraph{MTL, BL, G\"odel and Product fuzzy logics.}
\label{entry:p14-directory-fuzzy}
These systems use the exact implication, strong conjunction, lattice
connectives and constants of their linked Hilbert sheets. Their standard
consequence preserves value one. MTL ranges over all left-continuous
t-norms; BL over all continuous t-norms. The G\"odel standard chain uses
minimum as strong conjunction, while Product uses multiplication.
\Cref{def:p08-tnorm-residuum,tab:p08-tnorm-logics,
thm:p08-tnorm-completeness} give the operations, model classes and
premise boundaries. The stated arbitrary-premise standard completeness
for MTL and G\"odel logic is not transferred to BL, Product or the
standard infinite \L{}ukasiewicz chain. A fuzzy degree is not thereby
an additive probability \cite{Hajek1998,EstevaGodo2001,JeneiMontagna2002,
CignoliEstevaGodoTorrens2000,CintulaHajekNoguera2011}.

\paragraph{Supervaluation semantics.}
\label{entry:p14-directory-supervaluation}
Fine's specification spaces use complete admissible extensions of a
base specification. Supertruth is truth in every such completion; it
is not a third truth-functional value. The assumption that the base
point has an admissible completion matters. The exact semantics and
classical-tautology boundary are in
\cref{def:p08-supervaluation,thm:p08-supervaluation-classical}.
The directory does not identify truth preservation at each completion
with a different global preservation-of-supertruth relation; that
choice has to be declared before a consequence claim is made \cite{Fine1975}.

\section{Relevant, substructural, conditional, and defeasible systems}
\label{sec:p14-directory-relevant-defeasible}

\paragraph{Relevant implication E, R and RM.}
\label{entry:p14-directory-relevance}
The constant-free language has negation, conjunction, disjunction and
relevant implication. The matched pointed Routley--Meyer structures use
a ternary relation, a distinguished point, star and hereditary valuations.
The modular conditions for \(\mathbf E\), \(\mathbf R\) and
\(\mathbf{RM}\) are in \cref{def:p08-pointed-rm-frame}; the systems
have separate \hyperref[sheet:relevance-e]{E},
\hyperref[sheet:relevance-r]{R} and
\hyperref[sheet:relevance-rm]{RM} sheets. The stated adequacy is formula
theoremhood at the distinguished point, not an arbitrary-premise
local-consequence theorem on every broad ternary frame.
Variable sharing applies to E/R theorem implications under its
constant-free hypothesis; RM has the explicitly weaker alternative.
Full E and R theoremhood are undecidable, their pure implicational
fragments are decidable, and the separate RM decision result is cited.
Use \cref{entry:p08-relevance-r,thm:p08-relevance-adequacy,
thm:p08-relevance-undecidable,thm:p08-rm-decidable}.
Ackermann's gamma is a theoremhood metarule, not unrestricted disjunctive
syllogism from assumptions \cite{AndersonBelnap1975,RoutleyMeyer1973,
MeyerDunn1969,DunnRestall2002,Urquhart1984,Avron2016}.

\paragraph{IMLL, full classical linear logic and the affine variant.}
\label{entry:p14-directory-linear}
IMLL has atoms, tensor, tensor unit and linear implication; its finite
multiset contexts include exchange but not unrestricted weakening or
contraction. Its algebraic sequent semantics uses commutative residuated
partially ordered monoids. The exact adequacy/cut-admissibility treatment
retains its sketch boundary at \cref{entry:p08-imll}; the NP-complete
decision result is separately cited for that exact core.
Full one-sided CLL instead has multiplicatives, additives, their units,
paired atoms and exponentials. It uses formula-level linear negation
and the complete phase interpretation at
\cref{sec:p08-cll,def:p08-phase-space}. Its cited full adequacy and cut
elimination do not descend from the IMLL sketch. The full exponential
language, the multiplicative-additive fragment and the multiplicative
fragment have different decision classifications.
The commutative affine IMLL extension in
\cref{def:p08-affine-family-boundary} adds unrestricted weakening while
keeping contraction absent. Its sheet includes Cut; the treatment asserts
no independent affine adequacy or cut-elimination theorem
\cite{Girard1987,Girard1995,Lincoln1995,LincolnMitchellScedrovShankar1992}.

\paragraph{Associative Lambek calculus L.}
\label{entry:p14-directory-lambek}
The unit-free language has product and left and right divisions;
antecedents are nonempty ordered sequences. Residuated semigroups supply
the semantics: the ordered antecedent product is below the succedent
under every valuation. The treatment and cited adequacy, cut elimination
and NP-completeness are at
\cref{def:p08-lambek-semantics,entry:p08-lambek-l,
thm:p08-lambek-adequacy}. Exchange, weakening, contraction, the empty
antecedent and a unit are absent. The systems allowing empty antecedents
or adding a unit are distinct, not alternate labels for this card
\cite{Lambek1958,Buszkowski2010,Pentus2006}.

\paragraph{Basic conditional CK, unique selection and Lewis spheres.}
\label{entry:p14-directory-conditionals}
CK adds a primitive binary conditional to classical propositional logic.
Its set-valued selection model tests the consequent throughout the
selected set. Basic CK does not require that this set satisfy the
antecedent; adding that condition validates conditional identity and
changes the system. The equivalent principal-premise presentation and
cited completeness are in
\cref{def:p08-ck-semantics,entry:p08-conditional-ck}.
Stalnaker's separately specified selection clause picks one world for
each nonempty antecedent proposition and treats impossible antecedents
vacuously; centering is an extra condition. Lewis's nested spheres instead
use the printed variably strict, limit-free clause. V, VW and VC distinguish
no centering, weak centering and centering. The cited VC calculus and
adequacy are at \cref{entry:p08-counterfactual-vc}; its alternative name
C1 is not da Costa's \(C_1\). Unique selection, conditional excluded
middle and a least permitting sphere are not smuggled into VC
\cite{Chellas1975,Stalnaker1968,Lewis1971,Lewis1973,Nute1980,Wen2019}.

\paragraph{Standard deontic logic and other normative formalisms.}
\label{entry:p14-directory-deontic}
SDL is the monadic normal KD presentation with obligation read as box
over serial ideal alternatives and permission as its dual. Its formal
local consequence uses that model class; the reading of the operator
does not establish that every normative argument has this form.
\Cref{def:p08-sdl} and the following contrary-to-duty analysis retain
the Ross and Chisholm boundaries. Dyadic obligation and norm-based
input/output closure are separately described in
\cref{sec:p08-conditionals-deontic-free}; they are not SDL aliases or
truth-functional replacements for its monadic modality
\cite{VonWright1951,Chisholm1963,Hansson1969,McNamara2006}.

\paragraph{da Costa's propositional C1.}
\label{entry:p14-directory-da-costa}
This exact Hilbert language has negation, conjunction, disjunction and
implication, with a formula-level well-behavedness abbreviation.
Consequence is derivability on
\hyperref[sheet:hilbert-da-costa-c1]{the named C1 sheet}; the current
entry is not a claimed matrix adequacy theorem.
\Cref{def:p08-da-costa-c1,thm:p08-da-costa-c1-nonexplosion}
record nonexplosion and the extra consistency premise that licenses
controlled explosion. Its mechanism is different from LP's designated
values. Paraconsistency is a property of consequence, whereas
dialetheism is a claim about true contradictions \cite{daCosta1974,Priest2006}.

\paragraph{Preferential System P and rational consequence.}
\label{entry:p14-directory-preferential}
These relations connect classical antecedent and consequent formulas,
but test truth at the minimal antecedent states of a declared preference
model. Smoothness is part of the System P representation theorem.
Ranked models provide the stronger rational-monotony regime, not a
property of every smooth preference. The models, principles and cited
representation results are in
\cref{def:p08-preferential-model,thm:p08-klm-representation,
thm:p08-rational-representation}. Defeasible antecedent strengthening
may defeat a conclusion, so the inference is not monotonic classical
consequence \cite{KrausLehmannMagidor1990,LehmannMagidor1992}.

\paragraph{Reiter default logic, circumscription and autoepistemic logic.}
\label{entry:p14-directory-defaults}
Reiter theories combine classical premises with defaults whose prerequisites
must arise in grounded stages and whose justifications are tested against
the candidate extension. An extension is the stipulated grounded fixed
point; credulous and skeptical conclusions quantify over some and all
extensions respectively. Predicate circumscription instead selects models
minimal in declared predicates while keeping specified symbols fixed and
allowing others to vary. Moore's autoepistemic language adds belief and
uses a stable expansion fixed point with positive and negative introspective
assumptions. Its classical closure adds no KD45 axioms automatically.
Their actual interfaces and separately cited finite propositional bounds
are at \cref{entry:p08-reiter-default,entry:p08-circumscription,
entry:p08-autoepistemic}. None of the three operators is ordinary material
implication, and a failure of extension existence is not repaired by choosing
an unstated policy \cite{Reiter1980,McCarthy1980,Moore1985,Gottlob1992}.

\paragraph{AGM belief change.}
\label{entry:p14-directory-agm}
Revision and contraction operate on a deductively closed classical belief
set with an explicit formula input. A conclusion belongs to the resulting
belief set; there is no new object-language truth clause supplied solely
by the eight postulates. The postulates, Levi/Harper identities and finite
faithful-preorder representation are in
\cref{def:p08-agm,thm:p08-agm-representation}. Belief revision is therefore
located here as a formal change framework adjacent to consequence, not
identified with KLM preference or a source-independent decision algorithm
\cite{AlchourronGardenforsMakinson1985,KatsunoMendelzon1991}.

\section{Specialized formalisms for questions, programs, knowledge, and uncertainty}
\label{sec:p14-directory-specialized}

\paragraph{Questions: InqB and PMCE.}
\label{entry:p14-directory-questions}
InqB has a propositional language with inquisitive disjunction and support
semantics over information states. Implication quantifies over substates;
ordinary statement truth and question resolution are distinguished in
\cref{def:p08-inq-language,def:p08-inq-support,
sec:p08-inquisitive-propositional}. Its natural deduction is complete for
support consequence, and the finite-atom expressive/decision treatment
retains the restriction on legitimate substitutions for statement atoms.
PMCE instead has classical declaratives and separate ordered question
expressions with at least two syntactically distinct direct answers.
Evocation requires both coverage and unsettledness: the premises guarantee
some direct answer, but entail none individually. Its finite-premise calculus
is sound and complete for that relation at
\cref{sec:p18-question-evocation,thm:p18-qev-completeness}.
An added premise may settle an answer and destroy evocation. PMCE is not
InqB question-to-question support consequence, nor a guarantee that the
evoked question is worth investigating
\cite{Ciardelli2023InquisitivePropositional,Wisniewski2016QuestionEvocation}.

\paragraph{Temporal and fixed-point logics.}
\label{entry:p14-directory-temporal}
LTL uses next and strong until over positions of discrete infinite words.
Validity ranges over every word and position; its local finite-premise
interface and weak completeness are in
\cref{entry:p08-ltl,sec:p08-linear-temporal-completeness}.
CTL uses state formulas with each temporal constructor paired with a
path quantifier on total transition structures. CTL* separates state
and path syntax and admits general temporal path formulas beneath a
path quantifier. Their exact state/path interfaces and weak-adequacy
boundaries are at \cref{entry:p08-ctl,entry:p08-ctl-star}.
The modal mu-calculus uses positive fixed-point variables interpreted
as state sets, with least and greatest fixed points in the transition
model; its equational adequacy is not a theorem about arbitrary nonpositive
binders. Use \cref{def:p08-modal-mu-semantics,entry:p08-modal-mu}
and the distinct source constructions in
\cref{sec:p09-logic-in-computer-science}.
Finite traces, partial transition structures, different linear orders
and different alternation-depth conventions need their own clauses.
Kamp's strict since/until language and its order hypothesis are at
\cref{def:p08-tense-until-since,thm:p08-kamp}, not a license to extend
discrete LTL completeness to every order
\cite{GabbayPnueliShelahStavi1980,SistlaClarke1985,EmersonHalpern1986,
Reynolds2001,Kozen1983,Walukiewicz2000,Kamp1968}.

\paragraph{Knowledge, belief, common knowledge and announcements.}
\label{entry:p14-directory-epistemic}
Multi-agent S5 assigns an equivalence relation to each agent; multi-agent
KD45 instead uses serial, transitive, Euclidean relations for belief.
Local consequence preserves truth at the same world in every such
model. Their exact cards are \cref{entry:p08-s5n,entry:p08-kd45n};
the inherited S5 adequacy retains its sketch boundary.
Everyone-knows is a finite conjunction; distributed knowledge uses the
intersection of agent relations; common knowledge uses the
reflexive-transitive closure of their union. The latter is not a fixed
finite iteration of everyone-knows, and its finitary calculus is weakly
complete, not strongly complete for arbitrary premises when the group has
at least two agents; see \cref{entry:p08-common-knowledge,
def:p08-distributed-knowledge,thm:p08-common-knowledge-noncompact}.
PAL adds truthful public-announcement updates by restricting the model
to worlds satisfying the precondition. A false precondition makes its
box vacuously true. Its cited reduction to the exact static S5 base and
unsuccessful-update example are in \cref{entry:p08-pal}.
The reduction does not cover all action/event-model dynamic epistemic
logics \cite{Hintikka1962,FaginHalpernMosesVardi1995,Plaza1989,
BaltagMossSolecki1998,VanDitmarschVanDerHoekKooi2007}.

\paragraph{Hybrid, quantified-modal and free-logical interfaces.}
\label{entry:p14-directory-domain-variants}
Basic hybrid K(@) adds nominals naming one world and satisfaction
operators evaluating at that named world. Its current tableau excludes
binders and the global modality; see \cref{entry:p08-hybrid-k}.
Quantified modal logic is a parameterized family, not one entry obtained
by appending quantifiers to every modal card. The displayed Fitting
comparison fixes outer variable assignments, world-relative domains,
actualist relativization and the distinction between rigid denotation
and existence at a world. Its Barcan comparisons are at
\cref{def:p08-qml-family,def:p08-qml-model}; they do not specify a
family-wide complete calculus.
The negative, positive and neutral free-logic rows have different
atomic, identity and partial-term policies. Quantifiers range over their
inner existing domain, and consequence uses the row's designated values.
The exact variant-indexed relation and existence guards are at
\cref{def:p08-free-logic-consequence,tab:p08-free-logic-variants}.
Russellian contextual descriptions and primitive partial descriptions
are separately specified; a true predication in positive free logic
does not automatically establish inner-domain existence
\cite{ArecesTenCate2007,Fitting2002Alethic,FittingMendelsohn1998,
Lambert1967,Bencivenga1986,Lehmann2002,Russell1905}.

\paragraph{Description logics and finite-structure formalisms.}
\label{entry:p14-directory-description}
ALC uses Boolean concepts and existential/universal role restrictions,
interpreted as subsets and binary relations in a nonempty domain.
Knowledge-base consequence quantifies over interpretations satisfying
the actual TBox and ABox; standalone concept satisfiability supplies
neither component. The definition, terminating completion procedure,
general-TBox result and multimodal comparison are at
\cref{entry:p09-alc,def:p09-alc,thm:p09-alc-modal}.
The profile-specific EL, DL-Lite and OWL comparisons in
\cref{sec:p09-presentation-n-c-010} do not identify every profile with
ALC or transfer the standalone bound to general global assumptions.
Finite-model first-order logic, existential second-order logic and
fixed-point/query formalisms retain the finite relational vocabulary,
order and input conventions of \cref{sec:p09-finite-structures}.
An expressive capture theorem classifies properties of those structures,
not unrestricted logical validity \cite{Schild1991,SchmidtSchaussSmolka1991,
BaaderHorrocksLutzSattler2017,Libkin2004}.

\paragraph{Equational and contextual categorical logic.}
\label{entry:p14-directory-algebraic-categorical}
Equational logic uses finite terms of a finitary signature and equations
interpreted as identities in nonempty algebras. Each premise holds under
every valuation; this is not local implication at one fixed valuation.
The exact substitution/replacement calculus and its completeness proof are
at \cref{def:p10-equational-logic,thm:p10-equational-completeness}.
Pitts's contextual intuitionistic first-order presentation instead has
sorted terms, formula abstractions, term-equation judgments and formula
sequents with separate variable contexts and hypothesis lists.
Its categorical doctrine interpretation and soundness are at
\cref{def:p10-contextual-first-order,thm:p10-doctrine-soundness}; the
adjoint calculus is a separately matched presentation. A formula asserting
equality is not reflected automatically into equality of term arrows.
Boolean, Heyting, modal, MV and residuated algebras supply the
system-specific interpretations in \cref{sec:p10-nonclassical-algebras,
sec:p10-many-substructural}; the word ``categorical'' alone does not select
one consequence relation or make an arbitrary internal logic classical
\cite{BurrisSankappanavar2012,Pitts1995Draft,LambekScott1986}.

\paragraph{PDL, Hoare logic and separation logic.}
\label{entry:p14-directory-program}
PDL forms regular programs from atomic programs by composition, choice,
test and star, and evaluates program boxes over their relational successors.
\Cref{entry:p09-pdl,def:p09-pdl} fixes the test-containing language.
The later test-containing completeness source is not the earlier
test-free Kozen--Parikh presentation. Hoare triples instead relate
assertions to executions of the specified while language; partial
correctness and total correctness are distinct relations.
\Cref{def:p09-hoare-triple,thm:p09-hoare-soundness,
thm:p09-hoare-relative-completeness} retains soundness and the expressive
assertion-language conditions on relative completeness. Separation
logic uses heap splitting for separating conjunction and its separate
frame/locality conditions at \cref{def:p09-separation-logic,
thm:p09-frame-rule}. A relational program modality, a correctness triple
and a heap assertion are not three interchangeable syntactic forms
\cite{HarelKozenTiuryn2000,KozenParikh1981,Hoare1969,Cook1978,
OHearnReynoldsYang2001}.

\paragraph{Logic programs, databases and argumentation.}
\label{entry:p14-directory-logic-programs}
Definite Horn programs have the least Herbrand-model interpretation and
immediate-consequence operator in \cref{def:p09-definite-program}.
SLD is a search calculus with computed answers; a declarative least-model
result does not promise that one depth-first execution will terminate.
Finite normal stable models instead use the candidate-dependent reduct;
skeptical and credulous atom conclusions choose all and some stable models.
The well-founded semantics is a different, partial approximation.
Use \cref{def:p09-stable-model,def:p09-well-founded,
sec:p09-logic-programming}. Function-free Datalog evaluates a least
fixed point on the supplied finite database; conjunctive-query containment
is an all-databases relation, not evaluation on one database.
Dung frameworks and structured argumentation are separately treated in
Part XIV. Their extensions select acceptable arguments, and source-specific
bridges to default/program semantics preserve only their declared
objects and hypotheses. Argument acceptability or a successful dialogue
does not establish factual truth \cite{VanEmdenKowalski1976,Lloyd1987,
GelfondLifschitz1988,VanGelderRossSchlipf1991,ChandraMerlin1977,Dung1995}.

\paragraph{Atemporal likelihood logic, PCTL and CSL.}
\label{entry:p14-directory-probability}
The likelihood language forms Boolean combinations of rational linear
inequalities in the probabilities of propositional event formulas.
Its structures supply a probability space, valuations and measurable
definable events; consequence preserves inequality truth over that class.
The exact \(\mathsf{AX}_{\mathrm{MEAS}}\) source axiomatization and
small-model bound are at \cref{def:p09-likelihood-logic,
thm:p09-likelihood-logic}. No independence premise is added.
PCTL instead forms probability-threshold state formulas about temporal
path events in finite Markov chains or MDPs. The MDP syntax explicitly
chooses scheduler minimum or maximum. CSL uses a finite CTMC with
continuous-time residence and interval conventions. Their clauses and
model-checking bounds are at \cref{def:p09-dtmc-mdp,def:p09-pctl,
def:p09-ctmc-csl,thm:p09-csl-model-checking}.
CSL approximations justify a threshold verdict only with certified error
and endpoint treatment; equality is not settled by a displayed decimal.
These logics are not themselves Bayesian updating rules, causal
identification methods or models of human uncertainty judgments
\cite{FaginHalpernMegiddo1990,Halpern2003,BaierKatoen2008,
BaierHaverkortHermannsKatoen2003}.

\section{Principal results, limitations, and full-treatment references}
\label{sec:p14-directory-results}

Use a result through its actual assumptions and proof boundary. The following
routes distinguish a proof supplied in the manuscript, an explicitly
sketched argument and a source citation. They do not certify every
source-specific full presentation merely because a topic has a destination.
Each system treatment above supplies its language and consequence interface;
these result links supply the additional hypothesis-bearing statement and
the stated proof or source. Rules and side conditions remain on the exact
\cref{sec:p14-rule-sheets} destinations. Translations and representations
are collected in \cref{sec:p14-expressiveness-and-representation}; decision
questions and encodings in \cref{sec:p14-decidability-and-complexity}.

\begin{longtable}{@{}p{.25\textwidth}p{.38\textwidth}p{.29\textwidth}@{}}
\caption{Principal-result routes with their scope and current statement standing.}
\label{tab:p14-directory-results}\\
\toprule
System or question & Hypothesis-bearing result destination & Standing and limit\\
\midrule
\endfirsthead
\toprule
System or question & Hypothesis-bearing result destination & Standing and limit\\
\midrule
\endhead
\bottomrule
\endfoot
Classical propositional adequacy &
\cref{sec:prop-metatheory-01,sec:prop-metatheory-03}. & Inherited
soundness, completeness and compactness proofs; the fixed classical
calculus and two-valued semantics.\\\addlinespace
First-order adequacy and compactness &
\cref{sec:fol-metatheory-03,sec:fol-metatheory-04}. & The declared first-order
structure/assignment convention, not finite models or second-order full
semantics.\\\addlinespace
Full versus matched Henkin logic &
\cref{thm:p02-sol-full-noncompact,thm:p02-sol-derived-adequacy,
def:p02-sol-native-boundary}. & Full noncompactness; proved
translation-defined matched adequacy under ambient choice; native-rule
completeness not asserted.\\\addlinespace
IPL & \cref{entry:p08-ipl,sec:p08-ipl-kripke-completeness}. & Complete
prime-theory/canonical and finite-filtration developments; not a blanket
theorem for every constructive extension.\\\addlinespace
K, KT, S4, S5 & \cref{thm:p08-modal-canonical-completeness}. & Proved
strong local completeness for the four exact sheets; necessitation
remains theorem-only.\\\addlinespace
KD, KTB, K4 & \cref{thm:p08-kd-strong-completeness,
thm:p08-ktb-strong-completeness,thm:p08-k4-strong-completeness}. & Proved
separate canonical extensions over the declared countable language.\\\addlinespace
Neighborhood source systems & \cref{sec:p08-n-b-006-03}. & Complete
source-bound canonical constructions and their declared local premise
convention; not every independent neighborhood source variant.\\\addlinespace
K3, LP, FDE & \cref{entry:p08-k3,entry:p08-lp,entry:p08-fde}. & Proved
finite-matrix adequacy and decision specializations; fixed
constant-free connective language.\\\addlinespace
Real-valued fuzzy systems & \cref{thm:p08-tnorm-completeness}. & Cited
standard completeness with different finite/arbitrary-premise regimes;
all-algebra completeness is distinct.\\\addlinespace
E, R, RM & \cref{thm:p08-relevance-adequacy,
thm:p08-relevance-variable-sharing,thm:p08-rm-variable-sharing}. & Cited
theoremhood results on the matched pointed frames; no general
arbitrary-premise assertion.\\\addlinespace
IMLL & \cref{entry:p08-imll}. & Adequacy and cut admissibility
sketched for the exact multiplicative core; its separate cited decision
classification is not a full proof upgrade.\\\addlinespace
Full CLL and Lambek L & \cref{thm:p08-phase-adequacy,
thm:p08-cll-cut-elimination,thm:p08-lambek-adequacy}. & Cited results for
the complete declared phase or semigroup regime, not the affine extension
or a changed antecedent convention.\\\addlinespace
CK and VC & \cref{thm:p08-ck-completeness,thm:p08-vc-completeness}. &
Cited adequacy for the separate set-selection and centered limit-free
sphere classes.\\\addlinespace
System P and rational consequence & \cref{thm:p08-klm-representation,
thm:p08-rational-representation}. & Cited representation with smoothness
and ranking distinguished.\\\addlinespace
InqB and PMCE & \cref{thm:p08-inq-completeness,
thm:p08-inq-substitution,thm:p18-qev-completeness}. & Proved support
completeness and restricted substitution; separately proved finite-premise
evocation completeness.\\\addlinespace
LTL & \cref{sec:p08-linear-temporal-completeness}. & Complete validity
and finite-state decision developments for infinite discrete words;
arbitrary-premise strong completeness is not asserted.\\\addlinespace
CTL*, common knowledge & \cref{entry:p08-ctl-star,
thm:p08-common-knowledge-noncompact}. & Cited weak completeness and
noncompactness boundary; state/path types and the multi-agent hypothesis
remain essential.\\\addlinespace
PAL and basic hybrid K(@) & \cref{thm:p08-pal-reduction,
thm:p08-hybrid-tableau}. & Cited reduction/tableau adequacy for the
exact base and constructor set; no general event-model or binder result.\\\addlinespace
ALC & \cref{thm:p09-alc-tableau,thm:p09-alc-general-exptime}. & Proved
completion correctness/termination and general finite-input classification;
standalone concepts remain a distinct input.\\\addlinespace
PDL and Hoare logic & \cref{thm:p09-pdl-adequacy,
thm:p09-hoare-soundness,thm:p09-hoare-relative-completeness}. & Cited
test-containing PDL adequacy; proved Hoare soundness; cited relative
completeness conditional on the assertion language.\\\addlinespace
Likelihood and quantitative temporal logic & \cref{thm:p09-likelihood-logic,
thm:p09-pctl-model-checking,thm:p09-csl-model-checking,
thm:p09-csl-interval-computations}. & Cited likelihood axiomatization and
model-checking results; separately proved endpoint-aware CSL computation;
numerical threshold error remains explicit.\\
\end{longtable}

\paragraph{Selecting the full treatment.}
Follow the system entry for its semantic clauses and the named calculus for
its rules. Follow the result itself for proof status, and its annex link only
when that result assigns one. A cited result remains cited even when an
adjacent special case is proved. A source-bound alias is resolved within its
attribution: B/KB/KTB, neighborhood R/relevant R, Lewis C1/da Costa C1,
and full/Henkin second-order readings cannot be resolved by typography alone.
Countermodels and failed hypotheses are additional lookup routes in
\cref{sec:p14-countermodels-and-boundary-cases}; they do not replace the
definitions needed to interpret the failure.
